\documentclass[11pt,a4paper,twoside,openright]{report}
\usepackage{url,amsfonts,epsfig}
\usepackage{afterpage}
\usepackage{amsmath,amssymb}            
\usepackage{rotating}  
\usepackage{fancyhdr}  
\usepackage[scriptsize]{caption2} 
\hyphenation{a-gen-tiz-za-zio-ne}

\setlength{\paperwidth}{16cm}
\setlength{\paperheight}{24cm}
\setlength{\oddsidemargin} {2. cm}
\setlength{\evensidemargin} {2. cm}
\addtolength{\oddsidemargin} {-0.4 cm}
\addtolength{\evensidemargin} {-0.4 cm}
\linespread{1.1}
\usepackage[italian]{babel}
\usepackage[latin1]{inputenc}
\renewcommand{\captionfont}{\normalfont \sffamily \itshape \small}

\pagestyle{empty}



\begin{document}
\thispagestyle{empty}
%\begin{titlepage}
\vspace*{-1.5cm} 
\begin{center}
  \large
  UNIVERSIT\`A CA' FOSCARI VENEZIA\\
 \vspace*{1.3cm}


  FACOLT\`A DI SCIENZE \\
  MATEMATICHE, FISICHE E NATURALI\\
 \vspace*{0.3cm}
  \normalsize
  Dipartimento di Informatica\\

  \vspace*{0.3cm} \LARGE



  Tesi di Laurea in\\
 INFORMATICA\\
 \vspace*{.75truecm} \large
\bf{Mettere titolo}
\end{center}
\vspace*{3.0cm} \large
\begin{flushleft}

  Relatore: Dott.ssa Flaminia Luccio \\
  Correlatore: ... 

\end{flushleft}
\vspace*{1.0cm}
\begin{flushright}


  Tesi di Laurea di:\\ Nome, matricola numero \\ 
		       
\end{flushright}
\vspace*{0.5cm}
\begin{center}



  Anno Accademico 200N-200N+1
\end{center} \clearpage


\thispagestyle{empty} \normalfont \cleardoublepage
\vspace{17cm}

%\large
\begin{flushright}
\itshape{ Inserire eventuale dedica}
\end{flushright}

\thispagestyle{empty}  \cleardoublepage
\pagenumbering{Roman}
\chapter*{Ringraziamenti}

Ringrazio ................

\thispagestyle{empty} \vspace*{.75truecm} \normalfont \cleardoublepage
\pagestyle{plain}
\tableofcontents
\pagestyle{empty}\cleardoublepage
\pagenumbering{arabic}

\pagestyle{fancy} 
\fancyfoot{}                                               
\renewcommand{\chaptermark}[1]{\markboth{\chaptername\ \thechapter.\ #1}{}} 
\renewcommand{\sectionmark}[1]{\markright{\thesection.\ #1}}         
\fancyhead[LE,RO]{\bfseries\thepage}    
                                        
\fancyhead[RE]{\bfseries\leftmark}    
\fancyhead[LO]{\bfseries\rightmark}     
\renewcommand{\headrulewidth}{0.3pt} 

\newtheorem{problem}{{\sc Problema}}[chapter] 
\newtheorem{question}{{\sc Domanda}}[chapter]
\newtheorem{example}{{\sc Esempio}}[chapter]
\newtheorem{algo}{{\sc Algoritmo }} [chapter]
\newtheorem{step}{{\sc Passo }} [chapter]
\newtheorem{notaz}{{\sc Notazione}}[chapter]
\newtheorem{oss}{{\sc Osservazione}}[chapter]
\newtheorem{prop}{{\sc Proposizione}}[chapter]
\newtheorem{prope}{{\bf Propriet\`a}}[chapter]   
\newtheorem{definition}{Definizione}[chapter]
\newtheorem{teo}{Teorema}[chapter]
\newtheorem{lemm}{Lemma}[chapter]
\newtheorem{cor}{Corollario}[chapter]

%\def\qed{\ \ \vrule height6pt width5pt depth3pt}
\def\qed{\ \vrule height5pt width5pt depth0pt}
 
\newcommand{\pf}{{\sc Proof} }
\newcommand{\ep}{{\mbox{ }\nolinebreak{$\rule{2mm}{2mm}$}}}
\newcommand{\op}{\diamondsuit}
\newcommand{\cL}{{\cal L}}
\newcommand{\cN}{{\cal N}}
\newcommand{\cF}{{\cal F}}  
\newcommand{\QL}{Q_n [ {\cal L} ]}
\newcommand{\cNL}{{\cal N} [{\cal L} ]}
\newcommand{\diam}{{\rm diam}}
\newcommand{\pfo}{{\sc Proof} (outline) }
\newcommand{\OR}{\mbox{{\sc or}}}
%\newcommand{\TH}{\mbox{{\sc Th}}}
\newcommand{\iffe}{{\bf if}}
\newcommand{\then}{{\bf then}}
\newcommand{\els}{{\bf else}}
\newcommand{\beg}{{\bf begin}}
\newcommand{\e}{{\bf end}}   
\newcommand{\while}{{\bf while}}
\newcommand{\du}{{\bf do}}
\newcommand{\for}{{\bf for}}
\newcommand{\tu}{{\bf to}}   
\newcommand{\var}{{\bf var}} 
\newcommand{\of}{{\bf of}}   
\newcommand{\con}{{\bf with}}

\newcommand{\infig}[3]{\begin{figure*}[t]
\centerline{\epsfbox{#1}}
\caption{\label{#2} #3} \end{figure*}}







% si inseriscono i capitoli che sono inclusi in altri file

\include{capitolo1}
\include{capitolo2}
\include{capitolo3}

%il file comandi.tex contiene vari comandi latex.


\chapter{Introduzione}~\label{introduzione}
................

\section{Comandi vari}\label{comandi}


La sezione si richiama all'interno del testo con: Sezione~\ref{comandi}


\begin{itemize}
\item lista con *
\item ....
\end{itemize}


\begin{enumerate}
\item lista numerata
\item ....
\end{enumerate}

\begin{definition}
mettere una definizione
\end{definition}

\begin{teo}
mettere un teorema
\end{teo}
\pf la prova
\ep


Per fare tabelle e richiamarle Tabella~\ref{ds1}: 


\begin{table}
\baselineskip .6cm

\hrulefill

{\small

\begin{minipage}[c]{0.3\linewidth}

\begin{tabbing}
\\
scritta \\
scritta \\
\end{tabbing}
\end{minipage}

\hrulefill
\caption{ quello che si vuole scrivere.} \label{ds1}
}
\end{table}



per richiamare una figura nel testo ......Figure~\ref{etichettafigura}


\begin{figure}
 \centering \epsfig{file=nomefigura.eps, width=1.3in}
 \caption{scritta sotto la figura.\label{etichettafigura}}
 \end{figure}

Per allargare o restringere la figura cambiare width.

per creare la figura usare xfig \newline



{\bf scritta in neretto}, {\em scritta in italico}\newline


% commenta la riga \newline

Formule:


$1\leq i\geq l-1$  1<= i >= l-1






Per la bibligrafia seguire la traccia nell'ordine per riferimenti a un
libro, un lavoro incompleto a una 
conferenza e un giornale. Per richimarli nel testo usare \cite{H94,RCCR} :

\begin{thebibliography}{10}

\bibitem{H94}
F.~Harary. \newblock
Graph Theory. Reading, MA, Addison-Wesley, 1994.

\bibitem{L}   
T.~Lukka.  \newblock
A Gentle Introduction to Ted Nelson's ZigZag Structure (An incomplete
work in progress.), 2002 \newblock
{\tt http://www.nongnu.org/gzz/gi/gi.html}.

\bibitem{RCCR}

G.~Robertson, K.~Cameron, M.~Czerwinski and D.~Robbins, \newblock
Polyarchy visualization: visualizing multiple intersecting
hierarchies, \newblock
Proceedings of ACM CHI 2002 Conference on Human Factors in
Computing Systems, pp. 423--430, 2002.


\bibitem{N}
T.H.~Nelson. \newblock
A Cosmology for a different computer universe: data model~ mechanism,
virtual machine and visualization infrastructure.
Journal of Digital Information: Special Issue on Future Visions of
Common-Use Hypertext
anno
\end{thebibliography}


Se c'e' inserire l'enventuale appendice
\include{appendiceA}
\end{document}

